\section*{Chapitre \ref{ch:mx:benchmark-algorithms} --- Algorithmes à indice de référence}

\begin{solution}{exo:mx:benchmark-algorithms:1}
$15\,000\times5.50~\%=825$ actions dans la première tranche et $15\,000\times9.38~\%=1\,407$ dans la dernière.
\end{solution}

\begin{solution}{exo:mx:benchmark-algorithms:2}
Échanger $0.1\,o$ contre le $o$ des autres lui donne $0.1/1.1=9.09~\%$ du volume total. Pour être à 10~\% il faut échanger $r/(1-r)=11.1~\%$ du volume des autres
(ou 10~\% du volume total, le sien compris).
\end{solution}

\begin{solution}{exo:mx:benchmark-algorithms:3}
L'acheteur a payé 10.55 contre le VWAP de la journée de 10.75~: $-0.20$. Exécution~: $\tfrac12(10.10-10.00)+\tfrac12(11.00-11.00)=0.05$. Échéancier~:
$(0.5-0.25)\times10.00+(0.5-0.75)\times11.00=-0.25$~: il a trop acheté dans la tranche bon marché, une erreur de prévision heureuse.
\end{solution}

\begin{solution}{exo:mx:benchmark-algorithms:4}
La part exécution varie à peu près autant dans les deux types de journées (écarts-types de 0.2 à 0.24 ticks), tandis que la part échéancier est quatre fois plus grande
sur les jours de fort volume (2.17 contre 0.48 pour l'échéancier statique), car leurs pics remodèlent le volume de la journée. Plus la part échéancier est grande, plus
sa part dans une variance à laquelle la part exécution ajoute un montant fixe.
\end{solution}

\begin{solution}{exo:mx:benchmark-algorithms:5}
Son \omterm{def:mx:the-limit-order-book:orders}{ordre au marché} à la clôture a fait monter le cours de clôture de 12.6 ticks par rapport à la même journée sans cet ordre, si bien que l'indice de référence est monté jusqu'à ses exécutions~: les 5.9
ticks «~gagnés~» font partie de ce mouvement. Face au prix d'arrivée, qu'il ne peut pas faire bouger, il a payé 7.3 ticks, et son coût d'impact était de 6.8, le plus élevé.
\end{solution}

\begin{solution}{exo:mx:benchmark-algorithms:6}
Chacun échange $q=r(o+3q)$, donc $q/o=r/(1-3r)=0.2/0.4=50~\%$ du volume naturel, et le volume total vaut $o/(1-3r)=2.5$ fois le volume naturel. Pour
$r=25~\%$, $q/o=0.25/0.25=100~\%$ et le volume quadruple~; pour $r\ge1/3$ la boucle n'a pas de point fixe, et seule la taille des ordres l'arrête.
\end{solution}

\begin{solution}{exo:mx:benchmark-algorithms:7}
382 jours-action (10.1~\%) sont des jours de fort volume. Leurs tracking errors sont de 20.0 points de base pour la courbe statique et de 18.9 pour l'échéancier niveau plus AR
(17.7 et 17.2 avec des nouvelles sur 8~\% des jours)~: davantage de nouvelles rend les jours chargés plus fréquents et pires, et la mise à jour aide un peu plus.
\end{solution}

\begin{solution}{exo:mx:benchmark-algorithms:8}
Le VWAP inclut les propres transactions de l'algorithme, et dans le marché simulé ses achats ont fait monter le VWAP de la journée de 2.0 ticks~: son slippage se mesure
par rapport à un indice de référence qu'il a déplacé. Son coût d'impact, par rapport à la même journée sans lui, était de 2.7 ticks, et par rapport au prix d'arrivée il a payé 3.0.
\end{solution}

\begin{solution}{pb:mx:benchmark-algorithms:1}
\textbf{1.} Voir les définitions~; les indices de référence sont le VWAP de la journée, le prix pondéré par le temps, le VWAP sur la fenêtre d'exécution, le prix d'arrivée
et le cours de clôture.
\textbf{2.} Il a fait monter le VWAP de la journée de 2.0 ticks~; son coût d'impact était de 2.7 ticks contre 0.6 annoncé.
\textbf{3.} Son \omterm{def:mx:the-limit-order-book:orders}{ordre au marché} à la clôture a fait monter le cours de clôture de 12.6 ticks.
\textbf{4.} Le coût d'impact par rapport à la même journée sans l'algorithme (nombres aléatoires communs)~: il élimine les mouvements de prix propres à la journée, dont l'écart type
d'environ six ticks noie le coût mesuré par rapport à l'arrivée.
\textbf{5.} Vingt actions, 250 jours, 26 tranches~; profils en U, un facteur de marché, un choc au niveau de la journée, un écart spécifique AR(1), des nouvelles sur 8~\% des
jours-action~; un jour de fort volume a plus de 1.5 fois le volume moyen de l'action sur les 20 jours précédents (331 sur 3\,800 jours-action de test).
\textbf{6.} $\log v_k=\log(\bar Vp_k)+\ell+e_k$ avec $e$ AR(1)~; après chaque tranche le niveau est réestimé par moindres carrés généralisés et le dernier
écart s'estompe au taux $\rho$.
\textbf{7.} Sommation par parties~: $\sum(a_k-b_k)w_k=-\sum_{k<B}D_k(w_{k+1}-w_k)$.
\textbf{8.} 17.7 et 13.7 points de base pour la courbe statique, 17.2 et 13.3 pour le niveau plus AR.
\textbf{9.} Son AR(1) additif, ajusté surtout sur les jours calmes, traite un pic de nouvelles comme un bruit qui va revenir vers sa moyenne (20.6 contre 18.3 pour le PCA-ARMA statique).
\textbf{10.} $\sum\frac{q_k}{Q}(f_k-w_k)+\sum(\frac{q_k}{Q}-\frac{v_k}{V})w_k$.
\textbf{11.} Les seize premiers jours-action à fort volume de la période de test et seize autres répartis sur celle-ci~; le simulateur tire à l'avance les
ordres des \omterm{def:mx:why-there-is-a-spread:informed}{noise traders} à partir de la graine de la journée, de sorte que les deux échéanciers rencontrent le même flux.
\textbf{12.} Statique~: 2.24, 2.17 et 0.21 ticks~; dynamique~: 1.75, 1.71 et 0.20.
\textbf{13.} 0.66 et 0.73 ticks~: la demi-fourchette et l'impact que paie chaque \omterm{def:mx:the-almgren-chriss-framework:parent}{ordre enfant} au marché.
\textbf{14.} \emph{Résultat nommé}~: sur les jours de fort volume du simulateur, la part prévision de volume (échéancier) explique 99.1~\% de la variance du slippage du VWAP statique
et 98.6~\% de celle du \omterm{def:mx:benchmark-algorithms:dvwap}{VWAP dynamique} (écarts-types du slippage de 2.24 et 1.75 ticks)~; pour les autres jours, 73.9~\% et 68.5~\%.
\textbf{15.} Non~: l'échéancier dynamique a fait mieux sur 7 jours sur 16, la différence appariée du carré du slippage vaut 1.5 avec une erreur type de 1.8, et
une seule journée d'actualité porte l'essentiel du gain~; le panel de 3\,800 jours montre un gain réel mais faible (de 17.7 à 17.2 points de base).
\textbf{16.} $q=r(o+q)$ donne $q/o=r/(1-r)$.
\textbf{17.} Le rythme est passé de 21.1~\% à 24.6~\% du volume naturel et le coût d'impact de 6.4 à 10.0 ticks.
\textbf{18.} L'ordre a duré près de douze minutes au lieu de cinq, assez longtemps pour que la boucle atteigne son régime permanent (24.9~\%).
\textbf{19.} Que ses erreurs les jours chargés proviennent surtout de la prévision de volume, que la mise à jour en cours de journée aide un peu en moyenne et pas de façon fiable
sur une journée donnée, et qu'une meilleure prévision des jours d'actualité vaut plus qu'une mise à jour plus rapide.
\textbf{20.} De la qualité avec laquelle l'échéancier a suivi le volume de la journée et de la fourchette payée, et non de ce que l'ordre a coûté~: l'indice
de référence évolue avec l'ordre.
\end{solution}

\begin{solution}{iq:mx:benchmark-algorithms:1}
Un VWAP (ou un faible taux de participation) sur la journée~: pas d'\omterm{def:mx:the-almgren-chriss-framework:urgency}{urgence}, une action liquide, donc il faut minimiser l'impact et ressembler à la moyenne de la journée~; l'implementation
shortfall seulement avec une conviction qui s'effrite.
\iqlookfor{L'\omterm{def:mx:the-almgren-chriss-framework:urgency}{urgence} dicte le choix~; l'indice de référence suit.}
\end{solution}

\begin{solution}{iq:mx:benchmark-algorithms:2}
Partir du profil moyen de l'action sur les derniers jours, ajouter une estimation au niveau de la journée et une mise à jour intrajournalière (un modèle niveau plus AR, ou une composante commune
à toutes les actions plus une partie ARMA spécifique), avec un traitement particulier des événements connus (rééquilibrages d'indices, échéances). Le juger à l'aune de la
tracking error du VWAP qu'il produit, et non de l'erreur quadratique moyenne des volumes.
\iqlookfor{Un profil avec des mises à jour~; les événements~; une perte adaptée à l'usage.}
\end{solution}

\begin{solution}{iq:mx:benchmark-algorithms:3}
Sommation par parties~: la part échéancier vaut $-\sum D_k\Delta w_k$, donc une erreur dans une tranche corrigée à la suivante ne coûte que le mouvement de prix
entre les deux~; ce qui coûte, c'est être en avance ou en retard sur le volume pendant longtemps.
\iqlookfor{L'identité~; l'intuition d'être en avance ou en retard.}
\end{solution}

\begin{solution}{iq:mx:benchmark-algorithms:4}
Chaque algorithme compte les transactions de l'autre, si bien que les deux échangent plus vite que prévu ($r/(1-2r)$ du volume naturel) et supportent l'impact de l'autre. Plafonner le
taux de participation au volume des autres, hors l'autre grand acheteur s'il peut être identifié, abaisser ce taux, ou passer à un échéancier qui ne suit pas le volume.
\iqlookfor{La boucle de rétroaction~; une parade concrète.}
\end{solution}

\begin{solution}{iq:mx:benchmark-algorithms:5}
L'échéancier fournit des quantités cibles selon le temps (ou selon le volume), à partir de l'ordre, de l'état du marché et de ce qui a été exécuté~; la couche des \omterm{def:mx:the-almgren-chriss-framework:parent}{ordres enfants}
tente de les atteindre et renvoie les exécutions et le volume du marché. Il faut garder l'échéancier déterministe à entrées données, afin de pouvoir le rejouer.
\iqlookfor{Cibles et rétroaction~; testabilité.}
\end{solution}

\begin{solution}{iq:mx:benchmark-algorithms:6}
L'\omterm{def:mx:the-limit-order-book:orders}{ordre au marché} à la clôture a fait bouger le cours de clôture~: les exécutions collaient à l'indice de référence parce que celui-ci a monté jusqu'à elles, et le fonds, valorisé à cette
clôture, a acheté à un prix que son propre ordre avait gonflé~; le cours du lendemain peut retomber. Il faut regarder le déséquilibre publié avant la clôture et le
cours qui suit.
\iqlookfor{Contamination de l'indice de référence~; le déséquilibre de l'enchère de clôture.}
\end{solution}
